\subsection{Integração progressiva da solução}

A integração seguiu os marcos definidos no Capítulo~\ref{cap:metodo}, e a Tabela~\ref{tab:marcos_cronologia} registra quando cada um foi atingido. Os três primeiros marcos, que estabelecem a cadeia de ferramentas, a inicialização do transceptor e o enlace ponta a ponta, foram concluídos com a coleta de referência de 23 de agosto de 2026, com o quadro ainda em claro. O quinto e o sétimo, que acrescentam o envio ao \textit{backend} e o quadro autenticado, foram desenvolvidos em paralelo e validados juntos na coleta de 24 de setembro de 2026, descrita na Seção~\ref{sec:resultados_infra}. Sobre essa base, a frente de criptografia substitui o algoritmo da camada de cifra por uma opção de compilação, sem alterar o restante do firmware, e a frente de aprendizado de máquina consome a base persistida.

\begin{table}[htb]
\centering
\caption{Cronologia dos marcos da infraestrutura}
\label{tab:marcos_cronologia}
\begin{tabular}{clll}
\hline
\textbf{Marco} & \textbf{Objetivo} & \textbf{Data} & \textbf{Estado} \\
\hline
1 & Cadeia de ferramentas & jul. 2026 & concluído \\
2 & Rádio & jul. 2026 & concluído \\
3 & Enlace ponta a ponta & 23 ago. 2026 & concluído \\
4 & Sensores reais & & pendente \\
5 & Wi-Fi, HTTP e API & 24 set. 2026 & concluído \\
6 & Energia & & firmware pronto \\
7 & Segurança do enlace & 24 set. 2026 & Ascon pendente \\
\hline
\end{tabular}
\legend{Fonte: elaborado pelos autores (2026).}
\end{table}

Três dificuldades de integração merecem registro, porque afetam a coleta de dados. A primeira: abrir a porta serial do nó reinicia a placa, pois o conversor USB-serial pulsa as linhas de controle ao abrir a conexão; um reinício no meio da coleta zera a sessão. Desde então o nó opera em bateria e sem monitoramento durante as coletas, e o \textit{gateway} é o único ponto de observação. A segunda: o hotspot de celular usado como rede no laboratório é classificado pelo sistema operacional do servidor como rede pública, o que exige uma regra de firewall explícita para a porta da API. A terceira: o relógio do servidor estava cerca de meio segundo atrás do relógio do \textit{gateway}, sincronizado por NTP, o que impede usar a diferença entre o instante de recepção e o de ingestão como medida de latência sem sincronizar o servidor.


%%%%%%%%%%%%%%%%%%%%%%%%%%%%%%%%%%%%%%%%%%%%%%%%%%%%%%
%%%%%%%%%%%%%%%%%%%%%%%%%%%%%%%%%%%%%%%%%%%%%%%%%%%%%%
